\documentclass{noteformyself}

\newcommand{\Yang}[1]{\textcolor{red}{Yang: #1}}    % 批注命令
%一些常用的宏定义
\newcommand{\bbA}{{\mathbb{A}}}
\newcommand{\bba}{{\mathbb{a}}}
\newcommand{\bbB}{{\mathbb{B}}}
\newcommand{\bbb}{{\mathbb{b}}}
\newcommand{\bbC}{{\mathbb{C}}}
\newcommand{\bbc}{{\mathbb{c}}}
\newcommand{\bbD}{{\mathbb{D}}} 
\newcommand{\bbd}{{\mathbb{d}}}
\newcommand{\bbE}{{\mathbb{E}}}
\newcommand{\bbe}{{\mathbb{e}}}
\newcommand{\bbF}{{\mathbb{F}}}
\newcommand{\bbf}{{\mathbb{f}}}
\newcommand{\bbG}{{\mathbb{G}}}
\newcommand{\bbg}{{\mathbb{g}}}
\newcommand{\bbH}{{\mathbb{H}}}
\newcommand{\bbh}{{\mathbb{h}}}
\newcommand{\bbI}{{\mathbb{I}}}
\newcommand{\bbi}{{\mathbb{i}}}
\newcommand{\bbJ}{{\mathbb{J}}}
\newcommand{\bbj}{{\mathbb{j}}}
\newcommand{\bbK}{{\mathbb{K}}}
\newcommand{\bbk}{{\mathbb{k}}}
\newcommand{\bbL}{{\mathbb{L}}}
\newcommand{\bbl}{{\mathbb{l}}}
\newcommand{\bbM}{{\mathbb{M}}}
\newcommand{\bbm}{{\mathbb{m}}}
\newcommand{\bbN}{{\mathbb{N}}}
\newcommand{\bbn}{{\mathbb{n}}}
\newcommand{\bbO}{{\mathbb{O}}}
\newcommand{\bbo}{{\mathbb{o}}}
\newcommand{\bbP}{{\mathbb{P}}}
\newcommand{\bbp}{{\mathbb{p}}}
\newcommand{\bbQ}{{\mathbb{Q}}}
\newcommand{\bbq}{{\mathbb{q}}}
\newcommand{\bbR}{{\mathbb{R}}}
\newcommand{\bbr}{{\mathbb{r}}}
\newcommand{\bbS}{{\mathbb{S}}}
\newcommand{\bbs}{{\mathbb{s}}}
\newcommand{\bbT}{{\mathbb{T}}}
\newcommand{\bbt}{{\mathbb{t}}}
\newcommand{\bbU}{{\mathbb{U}}}
\newcommand{\bbu}{{\mathbb{u}}}
\newcommand{\bbV}{{\mathbb{V}}}
\newcommand{\bbv}{{\mathbb{v}}}
\newcommand{\bbW}{{\mathbb{W}}}
\newcommand{\bbw}{{\mathbb{w}}}
\newcommand{\bbX}{{\mathbb{X}}}
\newcommand{\bbx}{{\mathbb{x}}}
\newcommand{\bbY}{{\mathbb{Y}}}
\newcommand{\bby}{{\mathbb{y}}}
\newcommand{\bbZ}{{\mathbb{Z}}}
\newcommand{\bbz}{{\mathbb{z}}}
\newcommand{\bbone}{{\mathbb{1}}}

\newcommand{\calA}{{\mathcal{A}}}
\newcommand{\cala}{{\mathcal{a}}}
\newcommand{\calB}{{\mathcal{B}}}
\newcommand{\calb}{{\mathcal{b}}}
\newcommand{\calC}{{\mathcal{C}}}
\newcommand{\calc}{{\mathcal{c}}}
\newcommand{\calD}{{\mathcal{D}}}
\newcommand{\cald}{{\mathcal{d}}}
\newcommand{\calE}{{\mathcal{E}}}
\newcommand{\cale}{{\mathcal{e}}}
\newcommand{\calF}{{\mathcal{F}}}
\newcommand{\calf}{{\mathcal{f}}}
\newcommand{\calG}{{\mathcal{G}}}
\newcommand{\calg}{{\mathcal{g}}}
\newcommand{\calH}{{\mathcal{H}}}
\newcommand{\calh}{{\mathcal{h}}}
\newcommand{\calI}{{\mathcal{I}}}
\newcommand{\cali}{{\mathcal{i}}}
\newcommand{\calJ}{{\mathcal{J}}}
\newcommand{\calj}{{\mathcal{j}}}
\newcommand{\calK}{{\mathcal{K}}}
\newcommand{\calk}{{\mathcal{k}}}
\newcommand{\calL}{{\mathcal{L}}}
\newcommand{\call}{{\mathcal{l}}}
\newcommand{\calM}{{\mathcal{M}}}
\newcommand{\calm}{{\mathcal{m}}}
\newcommand{\calN}{{\mathcal{N}}}
\newcommand{\caln}{{\mathcal{n}}}
\newcommand{\calO}{{\mathcal{O}}}
\newcommand{\calo}{{\mathcal{o}}}
\newcommand{\calP}{{\mathcal{P}}}
\newcommand{\calp}{{\mathcal{p}}}
\newcommand{\calQ}{{\mathcal{Q}}}
\newcommand{\calq}{{\mathcal{q}}}
\newcommand{\calR}{{\mathcal{R}}}
\newcommand{\calr}{{\mathcal{r}}}
\newcommand{\calS}{{\mathcal{S}}}
\newcommand{\cals}{{\mathcal{s}}}
\newcommand{\calT}{{\mathcal{T}}}
\newcommand{\calt}{{\mathcal{t}}}
\newcommand{\calU}{{\mathcal{U}}}
\newcommand{\calu}{{\mathcal{u}}}
\newcommand{\calV}{{\mathcal{V}}}
\newcommand{\calv}{{\mathcal{v}}}
\newcommand{\calW}{{\mathcal{W}}}
\newcommand{\calw}{{\mathcal{w}}}
\newcommand{\calX}{{\mathcal{X}}}
\newcommand{\calx}{{\mathcal{x}}}
\newcommand{\calY}{{\mathcal{Y}}}
\newcommand{\caly}{{\mathcal{y}}}
\newcommand{\calZ}{{\mathcal{Z}}}
\newcommand{\calz}{{\mathcal{z}}}

\newcommand{\frakA}{{\mathfrak{A}}}
\newcommand{\fraka}{{\mathfrak{a}}}
\newcommand{\frakB}{{\mathfrak{B}}}
\newcommand{\frakb}{{\mathfrak{b}}}
\newcommand{\frakC}{{\mathfrak{C}}}
\newcommand{\frakc}{{\mathfrak{c}}}
\newcommand{\frakD}{{\mathfrak{D}}}
\newcommand{\frakd}{{\mathfrak{d}}}
\newcommand{\frakE}{{\mathfrak{E}}}
\newcommand{\frake}{{\mathfrak{e}}}
\newcommand{\frakF}{{\mathfrak{F}}}
\newcommand{\frakf}{{\mathfrak{f}}}
\newcommand{\frakG}{{\mathfrak{G}}}
\newcommand{\frakg}{{\mathfrak{g}}}
\newcommand{\frakH}{{\mathfrak{H}}}
\newcommand{\frakh}{{\mathfrak{h}}}
\newcommand{\frakI}{{\mathfrak{I}}}
\newcommand{\fraki}{{\mathfrak{i}}}
\newcommand{\frakJ}{{\mathfrak{J}}}
\newcommand{\frakj}{{\mathfrak{j}}}
\newcommand{\frakK}{{\mathfrak{K}}}
\newcommand{\frakk}{{\mathfrak{k}}}
\newcommand{\frakL}{{\mathfrak{L}}}
\newcommand{\frakl}{{\mathfrak{l}}}
\newcommand{\frakM}{{\mathfrak{M}}}
\newcommand{\frakm}{{\mathfrak{m}}}
\newcommand{\frakN}{{\mathfrak{N}}}
\newcommand{\frakn}{{\mathfrak{n}}}
\newcommand{\frakO}{{\mathfrak{O}}}
\newcommand{\frako}{{\mathfrak{o}}}
\newcommand{\frakP}{{\mathfrak{P}}}
\newcommand{\frakp}{{\mathfrak{p}}}
\newcommand{\frakQ}{{\mathfrak{Q}}}
\newcommand{\frakq}{{\mathfrak{q}}}
\newcommand{\frakR}{{\mathfrak{R}}}
\newcommand{\frakr}{{\mathfrak{r}}}
\newcommand{\frakS}{{\mathfrak{S}}}
\newcommand{\fraks}{{\mathfrak{s}}}
\newcommand{\frakT}{{\mathfrak{T}}}
\newcommand{\frakt}{{\mathfrak{t}}}
\newcommand{\frakU}{{\mathfrak{U}}}
\newcommand{\fraku}{{\mathfrak{u}}}
\newcommand{\frakV}{{\mathfrak{V}}}
\newcommand{\frakv}{{\mathfrak{v}}}
\newcommand{\frakW}{{\mathfrak{W}}}
\newcommand{\frakw}{{\mathfrak{w}}}
\newcommand{\frakX}{{\mathfrak{X}}}
\newcommand{\frakx}{{\mathfrak{x}}}
\newcommand{\frakY}{{\mathfrak{Y}}}
\newcommand{\fraky}{{\mathfrak{y}}}
\newcommand{\frakZ}{{\mathfrak{Z}}}
\newcommand{\frakz}{{\mathfrak{z}}}

\newcommand{\rmA}{{\mathrm{A}}}
\newcommand{\rma}{{\mathrm{a}}}
\newcommand{\rmB}{{\mathrm{B}}}
\newcommand{\rmb}{{\mathrm{b}}}
\newcommand{\rmC}{{\mathrm{C}}}
\newcommand{\rmc}{{\mathrm{c}}}
\newcommand{\rmD}{{\mathrm{D}}}
\newcommand{\rmd}{{\mathrm{d}}}
\newcommand{\rmE}{{\mathrm{E}}}
\newcommand{\rme}{{\mathrm{e}}}
\newcommand{\rmF}{{\mathrm{F}}}
\newcommand{\rmf}{{\mathrm{f}}}
\newcommand{\rmG}{{\mathrm{G}}}
\newcommand{\rmg}{{\mathrm{g}}}
\newcommand{\rmH}{{\mathrm{H}}}
\newcommand{\rmh}{{\mathrm{h}}}
\newcommand{\rmI}{{\mathrm{I}}}
\newcommand{\rmi}{{\mathrm{i}}}
\newcommand{\rmJ}{{\mathrm{J}}}
\newcommand{\rmj}{{\mathrm{j}}}
\newcommand{\rmK}{{\mathrm{K}}}
\newcommand{\rmk}{{\mathrm{k}}}
\newcommand{\rmL}{{\mathrm{L}}}
\newcommand{\rml}{{\mathrm{l}}}
\newcommand{\rmM}{{\mathrm{M}}}
\newcommand{\rmm}{{\mathrm{m}}}
\newcommand{\rmN}{{\mathrm{N}}}
\newcommand{\rmn}{{\mathrm{n}}}
\newcommand{\rmO}{{\mathrm{O}}}
\newcommand{\rmo}{{\mathrm{o}}}
\newcommand{\rmP}{{\mathrm{P}}}
\newcommand{\rmp}{{\mathrm{p}}}
\newcommand{\rmQ}{{\mathrm{Q}}}
\newcommand{\rmq}{{\mathrm{q}}}
\newcommand{\rmR}{{\mathrm{R}}}
\newcommand{\rmr}{{\mathrm{r}}}
\newcommand{\rmS}{{\mathrm{S}}}
\newcommand{\rms}{{\mathrm{s}}}
\newcommand{\rmT}{{\mathrm{T}}}
\newcommand{\rmt}{{\mathrm{t}}}
\newcommand{\rmU}{{\mathrm{U}}}
\newcommand{\rmu}{{\mathrm{u}}}
\newcommand{\rmV}{{\mathrm{V}}}
\newcommand{\rmv}{{\mathrm{v}}}
\newcommand{\rmW}{{\mathrm{W}}}
\newcommand{\rmw}{{\mathrm{w}}}
\newcommand{\rmX}{{\mathrm{X}}}
\newcommand{\rmx}{{\mathrm{x}}}
\newcommand{\rmY}{{\mathrm{Y}}}
\newcommand{\rmy}{{\mathrm{y}}}
\newcommand{\rmZ}{{\mathrm{Z}}}

\newcommand{\bfA}{{\mathbf{A}}}
\newcommand{\bfa}{{\mathbf{a}}}
\newcommand{\bfB}{{\mathbf{B}}}
\newcommand{\bfb}{{\mathbf{b}}}
\newcommand{\bfC}{{\mathbf{C}}}
\newcommand{\bfc}{{\mathbf{c}}}
\newcommand{\bfD}{{\mathbf{D}}}
\newcommand{\bfd}{{\mathbf{d}}}
\newcommand{\bfE}{{\mathbf{E}}}
\newcommand{\bfe}{{\mathbf{e}}}
\newcommand{\bfF}{{\mathbf{F}}}
\newcommand{\bff}{{\mathbf{f}}}
\newcommand{\bfG}{{\mathbf{G}}}
\newcommand{\bfg}{{\mathbf{g}}}
\newcommand{\bfH}{{\mathbf{H}}}
\newcommand{\bfh}{{\mathbf{h}}}
\newcommand{\bfI}{{\mathbf{I}}}
\newcommand{\bfi}{{\mathbf{i}}}
\newcommand{\bfJ}{{\mathbf{J}}}
\newcommand{\bfj}{{\mathbf{j}}}
\newcommand{\bfK}{{\mathbf{K}}}
\newcommand{\bfk}{{\mathbf{k}}}
\newcommand{\bfL}{{\mathbf{L}}}
\newcommand{\bfl}{{\mathbf{l}}}
\newcommand{\bfM}{{\mathbf{M}}}
\newcommand{\bfm}{{\mathbf{m}}}
\newcommand{\bfN}{{\mathbf{N}}}
\newcommand{\bfn}{{\mathbf{n}}}
\newcommand{\bfO}{{\mathbf{O}}}
\newcommand{\bfo}{{\mathbf{o}}}
\newcommand{\bfP}{{\mathbf{P}}}
\newcommand{\bfp}{{\mathbf{p}}}
\newcommand{\bfQ}{{\mathbf{Q}}}
\newcommand{\bfq}{{\mathbf{q}}}
\newcommand{\bfR}{{\mathbf{R}}}
\newcommand{\bfr}{{\mathbf{r}}}
\newcommand{\bfS}{{\mathbf{S}}}
\newcommand{\bfs}{{\mathbf{s}}}
\newcommand{\bfT}{{\mathbf{T}}}
\newcommand{\bft}{{\mathbf{t}}}
\newcommand{\bfU}{{\mathbf{U}}}
\newcommand{\bfu}{{\mathbf{u}}}
\newcommand{\bfV}{{\mathbf{V}}}
\newcommand{\bfv}{{\mathbf{v}}}
\newcommand{\bfW}{{\mathbf{W}}}
\newcommand{\bfw}{{\mathbf{w}}}
\newcommand{\bfX}{{\mathbf{X}}}
\newcommand{\bfx}{{\mathbf{x}}}
\newcommand{\bfY}{{\mathbf{Y}}}
\newcommand{\bfy}{{\mathbf{y}}}
\newcommand{\bfZ}{{\mathbf{Z}}}
\newcommand{\bfz}{{\mathbf{z}}}

\newcommand{\sfA}{{\mathsf{A}}}
\newcommand{\sfa}{{\mathsf{a}}}
\newcommand{\sfB}{{\mathsf{B}}}
\newcommand{\sfb}{{\mathsf{b}}}
\newcommand{\sfC}{{\mathsf{C}}}
\newcommand{\sfc}{{\mathsf{c}}}
\newcommand{\sfD}{{\mathsf{D}}}
\newcommand{\sfd}{{\mathsf{d}}}
\newcommand{\sfE}{{\mathsf{E}}}
\newcommand{\sfe}{{\mathsf{e}}}
\newcommand{\sfF}{{\mathsf{F}}}
\newcommand{\sff}{{\mathsf{f}}}
\newcommand{\sfG}{{\mathsf{G}}}
\newcommand{\sfg}{{\mathsf{g}}}
\newcommand{\sfH}{{\mathsf{H}}}
\newcommand{\sfh}{{\mathsf{h}}}
\newcommand{\sfI}{{\mathsf{I}}}
\newcommand{\sfi}{{\mathsf{i}}}
\newcommand{\sfJ}{{\mathsf{J}}}
\newcommand{\sfj}{{\mathsf{j}}}
\newcommand{\sfK}{{\mathsf{K}}}
\newcommand{\sfk}{{\mathsf{k}}}
\newcommand{\sfL}{{\mathsf{L}}}
\newcommand{\sfl}{{\mathsf{l}}}
\newcommand{\sfM}{{\mathsf{M}}}
\newcommand{\sfm}{{\mathsf{m}}}
\newcommand{\sfN}{{\mathsf{N}}}
\newcommand{\sfn}{{\mathsf{n}}}
\newcommand{\sfO}{{\mathsf{O}}}
\newcommand{\sfo}{{\mathsf{o}}}
\newcommand{\sfP}{{\mathsf{P}}}
\newcommand{\sfp}{{\mathsf{p}}}
\newcommand{\sfQ}{{\mathsf{Q}}}
\newcommand{\sfq}{{\mathsf{q}}}
\newcommand{\sfR}{{\mathsf{R}}}
\newcommand{\sfr}{{\mathsf{r}}}
\newcommand{\sfS}{{\mathsf{S}}}
\newcommand{\sfs}{{\mathsf{s}}}
\newcommand{\sfT}{{\mathsf{T}}}
\newcommand{\sft}{{\mathsf{t}}}
\newcommand{\sfU}{{\mathsf{U}}}
\newcommand{\sfu}{{\mathsf{u}}}
\newcommand{\sfV}{{\mathsf{V}}}
\newcommand{\sfv}{{\mathsf{v}}}
\newcommand{\sfW}{{\mathsf{W}}}
\newcommand{\sfw}{{\mathsf{w}}}
\newcommand{\sfX}{{\mathsf{X}}}
\newcommand{\sfx}{{\mathsf{x}}}
\newcommand{\sfY}{{\mathsf{Y}}}
\newcommand{\sfy}{{\mathsf{y}}}
\newcommand{\sfZ}{{\mathsf{Z}}}
\newcommand{\sfz}{{\mathsf{z}}}


\newcommand{\kk}{{\mathbf{k}}}
\newcommand{\kkk}{{\mathbb{k}}}
\newcommand{\KK}{{\mathbf{K}}}
\newcommand{\KKK}{{\mathbb{K}}}
\newcommand{\rkk}{{\kappa}} % residue field
\newcommand{\fkk}{{\mathscr{K}}} % fraction field

\renewcommand{\d}{{\mathrm{d}}}
\renewcommand{\i}{{\mathrm{i}}}
\renewcommand{\P}{\partial}


\newcommand{\Set}{\mathbf{Set}}
\newcommand{\Var}{\mathbf{Var}}
\newcommand{\Sch}{\mathbf{Sch}}
\newcommand{\Alg}{\mathbf{Alg}}
\newcommand{\Ring}{\mathbf{Ring}}
\newcommand{\Mod}{\mathbf{Mod}}
\newcommand{\TopMod}{\mathbf{TopMod}}
\newcommand{\Grp}{\mathbf{Grp}}
\newcommand{\Ab}{\mathbf{Ab}}
\newcommand{\Top}{\mathbf{Top}}
\newcommand{\Cat}{\mathbf{Cat}}
\newcommand{\AlgSp}{\mathbf{AlgSp}}
\newcommand{\Obj}{\operatorname{Obj}}


\newcommand{\ten}{\otimes}
\newcommand{\ratmap}{\dasharrow}
\newcommand{\injmap}{\hookrightarrow}
\newcommand{\surjmap}{\twoheadrightarrow}

\newcommand{\id}{{\operatorname{id}}}

\newcommand{\Frac}{\operatorname{Frac}}
\newcommand{\Der}{\operatorname{Der}}
\newcommand{\Spec}{\operatorname{Spec}}
\newcommand{\spec}{\operatorname{Spec}}
\newcommand{\mSpec}{\operatorname{mSpec}}
\newcommand{\depth}{\operatorname{depth}}
\newcommand{\idealht}{\operatorname{ht}}
\newcommand{\codim}{\operatorname{codim}}
\newcommand{\Supp}{\operatorname{Supp}}
\newcommand{\trdeg}{\operatorname{trdeg}}
\newcommand{\Ass}{\operatorname{Ass}}
\newcommand{\Ann}{\operatorname{Ann}}
\newcommand{\Ext}{\operatorname{Ext}}
\newcommand{\Tor}{\operatorname{Tor}}
\newcommand{\Hom}{\operatorname{Hom}}
\newcommand{\Eq}{\operatorname{Eq}}
\newcommand{\End}{\operatorname{End}}
\newcommand{\Mor}{\operatorname{Mor}}
\newcommand{\Mult}{\operatorname{Mult}}
\renewcommand{\Im}{\operatorname{Im}}
\newcommand{\Ker}{\operatorname{Ker}}
\newcommand{\rank}{\operatorname{rank}}
\newcommand{\cohdim}{\operatorname{coh.dim}}
\newcommand{\hldim}{\operatorname{hl.dim}}
\newcommand{\projdim}{\operatorname{proj.dim}}
\newcommand{\injdim}{\operatorname{inj.dim}}
\newcommand{\rad}{\operatorname{rad}}
\newcommand{\nil}{\operatorname{nil}}
\newcommand{\characteristic}{\operatorname{char}}
\newcommand{\Pic}{\operatorname{Pic}}
\newcommand{\NS}{\operatorname{NS}}
\newcommand{\Exc}{\operatorname{Exc}}
\newcommand{\Sing}{\operatorname{Sing}}
\newcommand{\Reg}{\operatorname{Reg}}
\newcommand{\Tr}{\operatorname{Tr}}
\newcommand{\Psef}{\operatorname{Psef}}
\newcommand{\Nef}{\operatorname{Nef}}
\newcommand{\Amp}{\operatorname{Amp}}
\newcommand{\Bigcone}{\operatorname{Big}}
\newcommand{\Frob}{\operatorname{Frob}}
\newcommand{\Bs}{\operatorname{Bs}}
\newcommand{\Stab}{\operatorname{Stab}}
\newcommand{\Sp}{\operatorname{Sp}}


\newcommand{\calHom}{{\mathcal{H\!o\!m}}}
\newcommand{\calExt}{{\mathcal{Ext}}}
\newcommand{\calProj}{{\mathcal{P\!r\!o\!j}}}
\newcommand{\calSpec}{{\mathcal{S\!p\!e\!c}}}


\newcommand{\mat}[4]{\left[ \begin{array}{cc}#1 &#2 \\ #3 &#4\end{array}\right]}
\newcommand{\threemat}[9]{\left[ \begin{array}{ccc}
  #1 & #2 & #3 \\
  #4 & #5 & #6 \\
  #7 & #8 & #9
\end{array}\right]}
\newcommand{\genmat}[9]{\left( \begin{array}{cccc}
  #1 & #2 & \cdots & #3 \\
  #4 & #5 & \cdots & #6 \\  
  \ldots & \ldots & \ddots & \ldots \\
  #7 & #8 & \cdots & #9
\end{array}\right)}





\newcommand{\contradiction}{
    \raisebox{-0.6ex}{\makebox[2.4ex][c]{\begin{tikzpicture}
        \draw (0,0) circle (1.2ex);
        \draw (0.7 ex, 0.7 ex) -- (-0.4 ex, 0.7 ex);
        \draw (-0.4 ex, 0.7 ex) -- (0.7 ex, -0.4 ex);
        \draw (0.7 ex, -0.4 ex) -- (0.7 ex, 0.7 ex);
        \draw (0.7 ex, 0.7 ex) -- (-0.848 ex, -0.848 ex);
    \end{tikzpicture}}}
    \ \ 
}

% legendre符号
\makeatletter
\def\legendre@dash#1#2{\hb@xt@#1{%
  \kern-#2\p@
  \cleaders\hbox{\kern.5\p@
    \vrule\@height.2\p@\@depth.2\p@\@width\p@
    \kern.5\p@}\hfil
  \kern-#2\p@
  }}
\def\@legendre#1#2#3#4#5{\mathopen{}\left(
  \sbox\z@{$\genfrac{}{}{0pt}{#1}{#3#4}{#3#5}$}%
  \dimen@=\wd\z@
  \kern-\p@\vcenter{\box0}\kern-\dimen@\vcenter{\legendre@dash\dimen@{#2}}\kern-\p@
  \right)\mathclose{}}
\newcommand\legendre[2]{\mathchoice
  {\@legendre{0}{1}{}{#1}{#2}}
  {\@legendre{1}{.5}{\vphantom{1}}{#1}{#2}}
  {\@legendre{2}{0}{\vphantom{1}}{#1}{#2}}
  {\@legendre{3}{0}{\vphantom{1}}{#1}{#2}}
}
\def\dlegendre{\@legendre{0}{1}{}}
\def\tlegendre{\@legendre{1}{0.5}{\vphantom{1}}}
\makeatother              % 加载符号定义文件
\addbibresource{../../Accessories/ref.bib}         % 加载参考文献文件


\title{Lecture of talk on Undergraduate Forum, ICCM25, Jan 07, 2026}
\author{Tianle Yang}
\date{\today}
% \authorpage{\href{www.tianleyang.com}{My Homepage}}

% % available on the mode "\sectionlevel=chapter" or "\sectionlevel=book".
% \setCJKfamilyfont{lxgwwenkai}{LXGW WenKai} % 定义霞鹜文楷，若未安装，请去掉相关代码编译或使用其他字体
% \coversentence{\CJKfamily{lxgwwenkai}阿巴阿巴阿巴阿巴阿巴阿巴阿巴阿巴!}
% \coverimage{placeholder.jpg} % 封面图片，
% \coverlinecolor{red} % 封面线条颜色，available on the mode "\sectionlevel=chapter" or "\sectionlevel=book".
% \covertitlefont{Allura} % 封面标题字体，available on the mode "\sectionlevel=chapter" or "\sectionlevel=book".
% \covertextcolor{green} % 封面文字颜色，available on the mode "\sectionlevel=chapter" or "\sectionlevel=book".


\def\acts{\ \rotatebox[origin=c]{-90}{$\circlearrowright$}\ }
\def\racts{\ \rotatebox[origin=c]{90}{$\circlearrowleft$}\ }
\def\dacts{\ \rotatebox[origin=c]{0}{$\circlearrowright$}\ }
\def\uacts{\ \rotatebox[origin=c]{180}{$\circlearrowright$}\ }


% \begin{document}

%     \maketitle

%     \section{Slide 1: Title}

%         Thanks for the introduction and the invitation of organizers.
%         It's my pleasure to give a talk here.

%         My talk is about about Algebraic Dynamics.

%     \bigskip

%     \section*{Slide 2: Kawaguchi--Silverman Conjecture}

%         Let \(X\) be a smooth projective variety defined over \(\mathbb{Q}\), and
%         let \(f : X \to X\) be a surjective endomorphism.
%         Fix an ample line bundle \(H\) on \(X\) and a height function \(h\)
%         associated to \(H\).

%         Kawaguchi and Silverman conjectured that if a point \(x \in X\) has dense orbit,
%         then the geometric growth rate of \(f\) equals the arithmetic growth rate
%         of the height along the orbit of \(x\).


%         The first one is the \emph{dynamical degree} $\delta_f$.
%         This is a geometric invariant, measuring how divisors grow under
%         iteration of $f$.

%         The second one is the \emph{arithmetic degree} $\alpha_f(x)$.
%         This measures how fast the height of a rational point grows along its
%         orbit.

%         The Kawaguchi--Silverman Conjecture can be viewed as an example of the slogan:

%         \begin{center}
%         \emph{Geometry controls arithmetic.}
%         \end{center}

%     \section*{Slide 3: An Example}

%         Let me illustrate this conjecture with a concrete example.

%         I drawn here the first few points in its orbit.

%         Since \(f^*H \sim 2H\), you can compute and see that the dynamical degree is \(2\).

%         And I list here the heights of the points in the orbit.

%         The ratio of the latter height to the former height is approximately \(2\),
%         so the arithmetic degree is also expected to be \(2\).

%     \section*{Slide 4: Known Cases}

%         The Kawaguchi--Silverman Conjecture is still open in general,
%         but it has been verified in many important settings.

%         In lower dimensions\dots
%         For some special classes of varieties\dots
%         And for some special classes of endomorphisms.

%     \section*{Slide 5: Main Results}

%         Now let me explain where our work fits into this picture.

%         We study the following special Fano fibration case, 
%         it comes from Minimal Model Program and can be viewed as a final output.

%         Explicitly, we consider a smooth Fano fibration
%         \[
%         \pi : X \to Y
%         \]
%         with relative Picard number one, where the base $Y$ is either an abelian
%         variety, or a smooth projective variety of Picard number one.

%         Under this setting, we show the KSC holds for $X$.

%         Since some sub cases are already known, we further impose some
%         dynamical conditions (dagger).

%         Under these assumptions,
%         we find a new equivariant fibration which preserves the dynamical degree.
%         And its base is polarized.
%         This will reduce the KSC for $X$ to the KSC for the base.
%         And the latter is known.

%     \section*{Slide 6: Sketch of Proof}

%         Let me very briefly explain the idea of the proof to the existence of this
%         new fibration.

%         We work in the special Fano fibration case, under the dynamical condition
%         that the dynamical degree strictly decreases.

%         First, under the dynamical restriction, we give a decomposition of the pseduo-effective cone of divisors on \(X\).
%         Use this decomposition and the geometry of base \(Y\), we show that the ramification divisor \(R_f\) is indeed an eigenvector of \(f^*\).
%         This shows that the induced dynamics on the base are polarized.

%         Next, we give a criterion for toric bundles by dynamics,
%         which applies to our Fano fibration setting.
%         We apply this criterion to show that the \(\kappa(X,R_f) > 0\).
%         This shows that the fibration we will construct is non-trivial.

%         Finally, using the Dynamical Iitaka theory, we combine these two steps to produce the desired fibration.

%     \section*{Slide 7: Conclusion}

%         That's all for my talk.
%         Thank you for your attention!

% \end{document}

\newcommand{\pausebreath}{\medskip}
\newcommand{\skipifshort}[1]{#1}

\begin{document}

    \maketitle

    \section{Slide 1: Title (25s)} % [20--25s]

        Thank you very much for the introduction and the invitation from the organizers.
        \pausebreath
        It is my pleasure to give a talk here.
        \pausebreath

        Today I will talk about some recent work in \emph{algebraic dynamics}.
        \pausebreath

    \section*{Slide 2: Kawaguchi--Silverman Conjecture (1min10s)} % [45--50s]

        Let \(X\) be a smooth projective variety defined over \(\mathbb{Q}\), and
        let \(f : X \to X\) be a surjective endomorphism.
        \pausebreath
        Fix an ample line bundle \(H\) on \(X\) and a height function \(h\)
        associated to \(H\).
        \pausebreath

        Kawaguchi and Silverman conjectured that if a point \(x \in X\) has a Zariski dense orbit,
        then the dynamical degree \(\delta_f\) coincides with the arithmetic degree \(\alpha_f(x)\).
        \pausebreath

        The first invariant is the \emph{dynamical degree} \(\delta_f\),
        which is the exponential growth rate of intersection numbers.
        \pausebreath

        The second one is the \emph{arithmetic degree} \(\alpha_f(x)\),
        which is the exponential growth rate of heights along the orbit of \(x\).
        \pausebreath

        % \skipifshort{
        % I will not explain these notions in detail today.
        % You can simply think of them as two growth rates: one geometric and one arithmetic.
        % \pausebreath
        % }

        The Kawaguchi--Silverman Conjecture can be viewed as an example of the slogan:
        \pausebreath

        \begin{center}
        \emph{Geometry controls arithmetic.}
        \end{center}

    \section*{Slide 3: An Example (2min)} % [45--50s]

        Let me illustrate this conjecture with a concrete example.
        \pausebreath

        On the slide, I have drawn the first few points in the orbit of a point.
        \pausebreath

        Since \(f^*H \sim 2H\), one can compute that the dynamical degree is equal to \(2\).
        \pausebreath

        I also list the heights of the points in the orbit.
        \pausebreath

        The ratio of successive heights is approximately \(2\),
        so the arithmetic degree is also expected to be \(2\),
        in agreement with the conjecture.
        \pausebreath

    \section*{Slide 4: Known Cases (2min25s)} % [20--25s]

        The Kawaguchi--Silverman Conjecture is still open in general,
        but it has been verified in many important cases.
        \pausebreath

        % \skipifshort{
        For example, it is known in lower dimensions,
        for certain special classes of varieties,
        and for certain special classes of endomorphisms.
        \pausebreath
        % }

        % I will not go into details here; this slide is only meant to show
        % that the conjecture has been widely studied.
        % \pausebreath

    \section*{Slide 5: Main Results (3min25s)} % [55--60s]

        Now let me explain where our work fits into this picture.
        \pausebreath

        We study a special Fano fibration case.
        This setting comes from the Minimal Model Program
        and can be viewed as one of its final outputs.
        \pausebreath

        Explicitly, we consider a smooth Fano fibration
        \[
        \pi : X \to Y
        \]
        with relative Picard number one,
        where the base \(Y\) is either an abelian variety,
        or a smooth projective variety of Picard number one.
        \pausebreath

        In this setting, we show that the Kawaguchi--Silverman Conjecture holds for \(X\).
        \pausebreath

        \skipifshort{
        Since some subcases are already known,
        we further impose certain dynamical conditions.
        \pausebreath
        }

        Under these assumptions, we construct a new equivariant fibration
        which preserves the dynamical degree,
        and whose base is polarized.
        \pausebreath

        This reduces the Kawaguchi--Silverman Conjecture for \(X\)
        to the conjecture for the base, which is already known.
        \pausebreath

    \section*{Slide 6: Sketch of Proof (4min25s)} % [45--50s]

        Let me very briefly explain the idea of the proof.
        \pausebreath

        We work in the special Fano fibration case,
        under the dynamical condition that the dynamical degree strictly decreases.
        \pausebreath

        First, using this dynamical restriction,
        we obtain a decomposition of the pseudo-effective cone of divisors on \(X\).
        \pausebreath

        Using this decomposition together with the geometry of the base \(Y\),
        we show that the ramification divisor \(R_f\) is an eigenvector of \(f^*\).
        \pausebreath

        This implies that the induced dynamics on the base are polarized.
        \pausebreath

        \skipifshort{
        Next, we establish a criterion for toric bundles using dynamical methods.
        \pausebreath
        }

        Applying this criterion, we show that \(\kappa(X, R_f) > 0\),
        which implies that the fibration we construct is non-trivial.
        \pausebreath

        Finally, using dynamical Iitaka theory,
        we combine these steps to produce the desired equivariant fibration.
        \pausebreath

    \section*{Slide 7: Conclusion (5min)} % [15--20s]

        To conclude, we use the dynamics to restrict geometry and get new structures.
        \pausebreath

        Thank you very much for your attention.

\end{document}
